\ifdefined\gkver\else\def\gkver{student}\fi
\documentclass[\gkver]{gaokaozhenti}
\newcommand{\ccnum}[1]{{\fontspec{Noto Sans CJK JP}#1}}
\begin{document}

\chapter{2015年全国课标Ⅱ}
%% paperType: 理综物理部分

\begin{enumerate}
\item
%% number: 14
%% paperName: 2015年全国课标Ⅱ第 14 题 6 分
%% typeId: 1
%% type: 单选题
%% chapter: 电场
%% point: 带电粒子的直线运动
%% method: 无
%% score: 6
%% degree: 650
%% duplicateId: 0
%% body:
如图，两平行的带电金属板水平放置。若在两板中间 a 点从静止释放一带电微粒，微粒恰好保持静止状态。现将两板绕过 a 点的轴（垂直于纸面）逆时针旋转 $45^\circ$，再由 a 点从静止释放一同样的微粒，该微粒将\xzanswer{D}

\onepicture[w=5.0cm]{figs/14.png}

\fourchoices[answer=D]
{保持静止状态}
{向左上方做匀加速运动}
{向正下方做匀加速运动}
{向左下方做匀加速运动}

%% answer: D
%% memo:
\memoanswer{
	本题原卷及材料中未提供详解，待补充。
}

\item
%% number: 15
%% paperName: 2015年全国课标Ⅱ第 15 题 6 分
%% typeId: 1
%% type: 单选题
%% chapter: 电磁感应
%% point: 切割磁感线电动势
%% method: 无
%% score: 6
%% degree: 450
%% duplicateId: 0
%% body:
如图，直角三角形金属框 abc 放置在匀强磁场中，磁感应强度大小为 $B$，方向平行于 ab 边向上。当金属框绕 ab 边以角速度 $\omega$ 逆时针转动时，a、b、c 三点的电势分别为 $\varphi_{a}$、$\varphi_{b}$、$\varphi_{c}$。已知 bc 边的长度为 $l$。下列判断正确的是\xzanswer{C}

\onepicture[h=4.0cm]{figs/15.png}

\fourchoices[answer=C]
{$\varphi_{a}>\varphi_{c}$，金属框中无电流}
{$\varphi_{b}>\varphi_{c}$，金属框中电流方向沿 a-b-c-a}
{$U_{bc}=-\frac{1}{2}Bl^{2}\omega$，金属框中无电流}
{$U_{ac}=\frac{1}{2}Bl^{2}\omega$，金属框中电流方向沿 a-c-b-a}

%% answer: C
%% memo:
\memoanswer{
	本题原卷及材料中未提供详解，待补充。
}

\item
%% number: 16
%% paperName: 2015年全国课标Ⅱ第 16 题 6 分
%% typeId: 1
%% type: 单选题
%% chapter: 曲线运动
%% point: 人造地球卫星
%% method: 无
%% score: 6
%% degree: 450
%% duplicateId: 0
%% body:
由于卫星的发射场不在赤道上，同步卫星发射后需要从转移轨道经过调整再进入地球同步轨道。当卫星在转移轨道上飞经赤道上空时，发动机点火，给卫星一附加速度，使卫星沿同步轨道运行。已知同步卫星的环绕速度约为 $3.1\times10^{3}\Ums$，某次发射卫星飞经赤道上空时的速度为 $1.55\times10^{3}\Ums$，此时卫星的高度与同步轨道的高度相同，转移轨道和同步轨道的夹角为 $30^\circ$，如图所示，发动机给卫星的附加速度的方向和大小约为\xzanswer{B}

\onepicture[w=7.0cm]{figs/16.png}

\fourchoices[answer=B]
{西偏北方向，$1.9\times10^{3}\Ums$}
{东偏南方向，$1.9\times10^{3}\Ums$}
{西偏北方向，$2.7\times10^{3}\Ums$}
{东偏南方向，$2.7\times10^{3}\Ums$}

%% answer: B
%% memo:
\memoanswer{
	本题原卷及材料中未提供详解，待补充。
}

\item
%% number: 17
%% paperName: 2015年全国课标Ⅱ第 17 题 6 分
%% typeId: 1
%% type: 单选题
%% chapter: 功和能
%% point: 交通工具启动
%% method: 无
%% score: 6
%% degree: 450
%% duplicateId: 0
%% body:
一汽车在平直公路上行驶。从某时刻开始计时，发动机的功率 $P$ 随时间 $t$ 的变化如图所示。假定汽车所受阻力的大小 $f$ 恒定不变。下列描述该汽车的速度 $v$ 随时间 $t$ 变化的图像中，可能正确的是\xzanswer{A}

\onepicture[h=3.2cm]{figs/17a.png}

\fourchoices[answer=A, ispicture=true, h=2.6cm]
{figs/17b.png}
{figs/17c.png}
{figs/17d.png}
{figs/17e.png}

%% answer: A
%% memo:
\memoanswer{
	本题原卷及材料中未提供详解，待补充。
}

\item
%% number: 18
%% paperName: 2015年全国课标Ⅱ第 18 题 6 分
%% typeId: 2
%% type: 多选题
%% chapter: 磁场
%% point: 电流的磁场
%% method: 无
%% score: 6
%% degree: 650
%% duplicateId: 0
%% body:
指南针是我国古代四大发明之一。关于指南针，下列说明正确的是\xzanswer{BC}

\fourchoices[answer=BC]
{指南针可以仅具有一个磁极}
{指南针能够指向南北，说明地球具有磁场}
{指南针的指向会受到附近铁块的干扰}
{在指南针正上方附近沿指针方向放置一直导线，导线通电时指南针不偏转}

%% answer: BC
%% memo:
\memoanswer{
	本题原卷及材料中未提供详解，待补充。
}

\item
%% number: 19
%% paperName: 2015年全国课标Ⅱ第 19 题 6 分
%% typeId: 2
%% type: 多选题
%% chapter: 磁场
%% point: 带电粒子在磁场中的运动
%% method: 无
%% score: 6
%% degree: 450
%% duplicateId: 0
%% body:
有两个运强磁场区域Ⅰ和Ⅱ，Ⅰ中的磁感应强度是Ⅱ中的 $k$ 倍，两个速率相同的电子分别在两磁场区域做圆周运动。与Ⅰ中运动的电子相比，Ⅱ中的电子\xzanswer{AC}

\fourchoices[answer=AC]
{运动轨迹的半径是Ⅰ中的 $k$ 倍}
{加速度的大小是Ⅰ中的 $k$ 倍}
{做圆周运动的周期是Ⅰ中的 $k$ 倍}
{做圆周运动的角速度是Ⅰ中的相等}

%% answer: AC
%% memo:
\memoanswer{
	本题原卷及材料中未提供详解，待补充。
}

\item
%% number: 20
%% paperName: 2015年全国课标Ⅱ第 20 题 6 分
%% typeId: 2
%% type: 多选题
%% chapter: 运动和力
%% point: 连接体
%% method: 无
%% score: 6
%% degree: 450
%% duplicateId: 0
%% body:
在一东西向的水平直铁轨上，停放着一列已用挂钩链接好的车厢。当机车在东边拉着这列车厢以大小为 $a$ 的加速度向东行驶时，链接某两相邻车厢的挂钩 P 和 Q 间的拉力大小为 $F$；当机车在西边拉着这列车厢以大小为 $\frac{2}{3}a$ 的加速度向东行驶时，P 和 Q 间的拉力大小仍为 $F$。不计车厢与铁轨间的摩擦，每节车厢质量相同，则这列车厢的节数可能为\xzanswer{BC}

\fourchoices[answer=BC]
{8}
{10}
{15}
{18}

%% answer: BC
%% memo:
\memoanswer{
	本题原卷及材料中未提供详解，待补充。
}

\item
%% number: 21
%% paperName: 2015年全国课标Ⅱ第 21 题 6 分
%% typeId: 2
%% type: 多选题
%% chapter: 功和能
%% point: 整体机械能守恒
%% method: 无
%% score: 6
%% degree: 250
%% duplicateId: 0
%% body:
如图，滑块 a、b 的质量均为 $m$，a 套在固定直杆上，与光滑水平地面相距 $h$，b 放在地面上，a、b 通过铰链用刚性轻杆连接。不计摩擦，a、b 可视为质点，重力加速度大小为 $g$。则\xzanswer{BD}

\onepicture[h=3.2cm]{figs/21.png}

\fourchoices[answer=BD]
{a 落地前，轻杆对 b 一直做正功}
{a 落地时速度大小为 $\sqrt{2gh}$}
{a 下落过程中，其加速度大小始终不大于 $g$}
{a 落地前，当 a 的机械能最小时，b 对地面的压力大小为 $mg$}

%% answer: BD
%% memo:
\memoanswer{
	本题原卷及材料中未提供详解，待补充。
}

\item
%% number: 22
%% paperName: 2015年全国课标Ⅱ第 22 题 6 分
%% typeId: 4
%% type: 实验题
%% chapter: 直线运动
%% point: 打点计时器公式
%% method: 无
%% score: 6
%% degree: 650
%% duplicateId: 0
%% body:
某学生用图（a）所示的实验装置测量物块与斜面的动摩擦因数。已知打点计时器所用电源的频率为 $50\UHz$，物块下滑过程中所得到的纸带的一部分如图（b）所示，图中标出了 5 个连续点之间的距离。

\begin{enumerate}
	\item
	物块下滑时的加速度 $a=$\tkanswer[1.8cm]{3.25}$\Umsq$，打点 C 点时物块的速度 $v=$\tkanswer[1.8cm]{1.79}$\Ums$；

	\item
	已知重力加速度大小为 $g$，为求出动摩擦因数，还需测量的物理量是\tkanswer[1.2cm]{C}（填正确答案标号）
	（A）物块的质量　（B）斜面的高度　（C）斜面的倾角
\end{enumerate}

\twopicture[showlabel=false, hA=3.2cm, hB=2.6cm]
{figs/22a.png}
{figs/22b.png}

%% answer:
\jdanswer{
（1）$3.25\Umsq$，$1.79\Ums$；（2）C
}
%% memo:
\memoanswer{
	本题原卷及材料中未提供详解，待补充。
}

\item
%% number: 23
%% paperName: 2015年全国课标Ⅱ第 23 题 9 分
%% typeId: 4
%% type: 实验题
%% chapter: 电路
%% point: 电路实验
%% method: 无
%% score: 9
%% degree: 450
%% duplicateId: 0
%% body:
电压表满偏时通过该表的电流是半偏时通过该表的电流的两倍。某同学利用这一事实测量电压表的内阻（半偏法），实验室提供材料器材如下：

待测电压表 V（量程 $3\UV$，内阻约为 $3000\UO$），电阻箱 $R_{0}$（最大阻值为 $99999.9\UO$），滑动变阻器 $R_{1}$（最大阻值 $100\UO$，额定电流 $2\UA$），电源 $E$（电动势 $6\UV$，内阻不计），开关 2 个，导线若干。

\begin{enumerate}
	\item
	虚线框内为该同学设计的测量电压表内阻的电路图的一部分，将电路图补充完整。

	\item
	根据设计的电路，写出实验步骤：\tkanswer[4.0cm]{见解析（2）}。

	\item
	将这种方法测出的电压表内阻记为 $R_{V}'$，与电压表内阻的真实值 $R_{V}$ 相比，$R_{V}'$\tkanswer[1.0cm]{$>$}$R_{V}$（填“$>$”、“$=$”或“$<$”），主要理由是\tkanswer[4.0cm]{见解析（3）}。
\end{enumerate}

\onepicture[align=right, h=4.6cm]{figs/23.png}

%% answer:
\jdanswer{
\begin{enumerate}
	\item
	见解析图。

	\item
	见解析。

	\item
	$>$
\end{enumerate}
}
%% memo:
\memoanswer{
	（1）答案电路图如下：

	\onepicture[h=4.6cm]{figs/23b.png}

	（2）移动滑动变阻器的滑片，以保证通电后电压表所在支路分压最小；闭合开关 $S_{1}$、$S_{2}$，调节 $R_{1}$，使电压表指针满偏；保证滑动变阻器的滑片的位置不变，断开 $S_{2}$，调节电阻箱 $R_{0}$，使电压表的指针半偏；读取电阻箱所示的电阻值，即为测得的电压表内阻。

	（3）断开 $S_{2}$，调节电阻箱使电压表成半偏状态，电压表所在支路的总电阻增大，分得的电压也增大，此时 $R_{0}$ 两端的电压大于电压表的半偏电压，故 $R_{V}'>R_{V}$。
}

\item
%% number: 24
%% paperName: 2015年全国课标Ⅱ第 24 题 12 分
%% typeId: 6
%% type: 计算题
%% chapter: 电场
%% point: 带电粒子的偏转
%% method: 无
%% score: 12
%% degree: 650
%% duplicateId: 0
%% body:
如图，一质量为 $m$、电荷量为 $q$（$q>0$）的粒子在匀强电场中运动，A、B 为其运动轨迹上的两点。已知该粒子在 A 点的速度大小为 $v_{0}$，方向与电场方向的夹角为 $60^\circ$；它运动到 B 点时速度方向与电场方向的夹角为 $30^\circ$。不计重力。求 A、B 两点间的电势差。

\onepicture[align=right, h=4.0cm]{figs/24.png}

%% answer:
\jdanswer{
$U_{AB}=\frac{mv_{0}^{2}}{q}$
}
%% memo:
\memoanswer{
	设带电粒子在 B 点的速度大小为 $v_{B}$，粒子在垂直于电场方向的分速度不变，即
	$v_{B}\sin30^\circ=v_{0}\sin60^\circ$，
	由此得 $v_{B}=\sqrt{3}v_{0}$。
	设 A、B 两点间的电势差为 $U_{AB}$，由动能定理有
	$qU_{AB}=\frac{1}{2}mv_{B}^{2}-\frac{1}{2}mv_{0}^{2}$，
	解得 $U_{AB}=\frac{mv_{0}^{2}}{q}$。
}

\item
%% number: 25
%% paperName: 2015年全国课标Ⅱ第 25 题 20 分
%% typeId: 6
%% type: 计算题
%% chapter: 运动和力
%% point: 牛顿定律的应用
%% method: 无
%% score: 20
%% degree: 250
%% duplicateId: 0
%% body:
下暴雨时，有时会发生山体滑坡或泥石流等地质灾害。某地有一倾角为 $\theta=37^\circ$（$\sin37^\circ=0.6$）的山坡 C，上面有一质量为 $m$ 的石板 B，其上下表面与斜坡平行；B 上有一碎石堆 A（含有大量泥土），A 和 B 均处于静止状态，如图所示。假设某次暴雨中，A 浸透雨水后总质量也为 $m$（可视为质量不变的滑块），在极短时间内，A、B 间的动摩擦因数 $\mu_{1}$ 减小为 $\frac{3}{8}$，B、C 间的动摩擦因数 $\mu_{2}$ 减小为 0.5，A、B 开始运动，此时刻为计时起点；在第 $2\Us$ 末，B 的上表面突然变为光滑，$\mu_{2}$ 保持不变。已知 A 开始运动时，A 离 B 下边缘的距离 $l=27\Um$，C 足够长，设最大静摩擦力等于滑动摩擦力。取重力加速度大小 $g=10\Umsq$。求：

\begin{enumerate}
	\item
	在 $0\sim2\Us$ 时间内 A 和 B 加速度的大小；

	\item
	A 在 B 上总的运动时间。
\end{enumerate}

\onepicture[align=right, h=2.6cm]{figs/25.png}

%% answer:
\jdanswer{
\begin{enumerate}
	\item
	$a_{1}=3\Umsq$，$a_{2}=1\Umsq$

	\item
	$t_{\text{总}}=4\Us$
\end{enumerate}
}
%% memo:
\memoanswer{
	（1）在 $0\sim2\Us$ 时间内，A 和 B 的受力如图所示，其中 $f_{1}$、$N_{1}$ 是 A 与 B 之间的摩擦力和正压力的大小，$f_{2}$、$N_{2}$ 是 B 与 C 之间的摩擦力和正压力的大小，方向如图所示。由滑动摩擦力公式和力的平衡条件

	$f_{1}=\mu_{1}N_{1}$\ccnum{①}，$N_{1}=mg\cos\theta$\ccnum{②}，$f_{2}=\mu_{2}N_{2}$\ccnum{③}，

	\onepicture[h=2.4cm]{figs/25b.png}

	$N_{2}=N_{1}+mg\cos\theta$\ccnum{④}，

	规定沿斜面向下为正。设 A 和 B 的加速度分别为 $a_{1}$ 和 $a_{2}$，由牛顿第二定律得

	$mg\sin\theta-f_{1}=ma_{1}$\ccnum{⑤}，$mg\sin\theta-f_{2}+f_{1}=ma_{2}$\ccnum{⑥}，

	联立\ccnum{①②③④⑤⑥}式，并代入题给条件得

	$a_{1}=3\Umsq$\ccnum{⑦}，$a_{2}=1\Umsq$\ccnum{⑧}。

	（2）在 $t_{1}=2\Us$ 时，设 A 和 B 的速度分别为 $v_{1}$ 和 $v_{2}$，则

	$v_{1}=a_{1}t_{1}=6\Ums$\ccnum{⑨}，$v_{2}=a_{2}t_{1}=2\Ums$\ccnum{⑩}。

	$t>t_{1}$ 时，设 A 和 B 的加速度分别为 $a_{1}'$ 和 $a_{2}'$。此时 A 与 B 之间摩擦力为零，同理可得

	$a_{1}'=6\Umsq$\ccnum{⑪}，$a_{2}'=-2\Umsq$\ccnum{⑫}，

	即 B 做减速运动。设经过时间 $t_{2}$，B 的速度减为零，则有

	$v_{2}+a_{2}'t_{2}=0$\ccnum{⑬}，

	联立\ccnum{⑩⑫⑬}式得

	$t_{2}=1\Us$\ccnum{⑭}，

	在 $t_{1}+t_{2}$ 时间内，A 相对于 B 运动的距离为

	$s=\left(\frac{1}{2}a_{1}t_{1}^{2}+v_{1}t_{2}+\frac{1}{2}a_{1}'t_{2}^{2}\right)-\left(\frac{1}{2}a_{2}t_{1}^{2}+v_{2}t_{2}+\frac{1}{2}a_{2}'t_{2}^{2}\right)$\ccnum{⑮}
	$=12\Um<27\Um$。

	此后 B 静止不动。A 继续在 B 上滑动。设再经过时间 $t_{3}$ 后 A 离开 B，则有

	$l-s=\left(v_{1}+a_{1}'t_{2}\right)t_{3}+\frac{1}{2}a_{1}'t_{3}^{2}$\ccnum{⑯}，

	可得 $t_{3}=1\Us$（另一解不合题意，舍去）\ccnum{⑰}。

	设 A 在 B 上总的运动时间为 $t_{\text{总}}$，有

	$t_{\text{总}}=t_{1}+t_{2}+t_{3}=4\Us$\ccnum{⑱}。

	（利用下面的速度图象求解，正确的，参照上述答案参考给分）

	\onepicture[h=4.2cm]{figs/25c.png}
}

\item
%% number: 33
%% paperName: 2015年全国课标Ⅱ第 33 题 10 分
%% typeId: 6
%% type: 计算题
%% chapter: 气体性质
%% point: 玻意耳定律
%% method: 无
%% score: 10
%% degree: 450
%% duplicateId: 0
%% body:
【物理——选修 3-3】（15 分）

（1）（5 分）关于扩散现象，下列说法正确的是\xzanswer{ACD}（填正确答案标号。选对 1 个得 2 分，选对 2 个得 4 分，选对三个得 5 分。每选错一个扣 3 分，最低得分为 0 分）

\fivechoices[answer=ACD]
{温度越高，扩散进行得越快}
{扩散现象是不同物质间的一种化学反应}
{扩散现象是由物质分子无规则运动产生的}
{扩散现象在气体、液体和固体中都能发生}
{液体中的扩散现象是由于液体的对流形成的}

（2）（10 分）如图，一粗细均匀的 U 形管竖直放置，A 侧上端封闭，B 侧上端与大气相通，下端开口处开关 K 关闭；A 侧空气柱的长度为 $l=10.0\Ucm$，B 侧水银面比 A 侧的高 $h=3.0\Ucm$。现将开关 K 打开，从 U 形管中放出部分水银，当两侧水银面的高度差为 $h_{1}=10.0\Ucm$ 时将开关 K 关闭。已知大气压强 $p_{0}=75.0\UcmHg$。求：

\begin{enumerate}
	\item
	求放出部分水银后 A 侧空气柱的长度；

	\item
	此后再向 B 侧注入水银，使 A、B 两侧的水银面达到同一高度，求注入的水银在管内的长度。
\end{enumerate}

\onepicture[h=5.0cm]{figs/33.png}

%% answer:
\jdanswer{
\begin{enumerate}
	\item
	$12.0\Ucm$

	\item
	$13.2\Ucm$
\end{enumerate}
}
%% memo:
\memoanswer{
	以 $\UcmHg$ 为压强单位。设 A 侧空气柱长度 $l=10.0\Ucm$ 时的压强为 $p$；当两侧水银面的高度差为 $h_{1}=10.0\Ucm$ 时，空气柱的长度为 $l_{1}$，压强为 $p_{1}$。由玻意耳定律得

	$pl=p_{1}l_{1}$\ccnum{①}，

	由力学平衡条件得

	$p=p_{0}+h$\ccnum{②}，

	打开开关 K 放出水银的过程中，B 侧水银面处的压强始终为 $p_{0}$，而 A 侧水银面处的压强随空气柱长度的增加逐渐减小，B、A 两侧水银面的高度差也随之减小，直至 B 侧水银面低于 A 侧水银面 $h_{1}$ 为止。由力学平衡条件有

	$p_{1}=p_{0}-h_{1}$\ccnum{③}，

	联立\ccnum{①②③}式，并代入题给数据得

	$l_{1}=12.0\Ucm$\ccnum{④}。

	（ii）当 A、B 两侧的水银面达到同一高度时，设 A 侧空气柱的长度为 $l_{2}$，压强为 $p_{2}$。由玻意耳定律得

	$pl=p_{2}l_{2}$\ccnum{⑤}，

	由力学平衡条件有

	$p_{2}=p_{0}$\ccnum{⑥}，

	联立\ccnum{②⑤⑥}式，并代入题给数据得

	$l_{2}=10.4\Ucm$\ccnum{⑦}，

	设注入的水银在管内的长度为 $\Delta h$，依题意得

	$\Delta h=2(l_{1}-l_{2})+h_{1}$\ccnum{⑧}，

	联立\ccnum{④⑦⑧}式，并代入题给数据得

	$\Delta h=13.2\Ucm$\ccnum{⑨}。
}

\item
%% number: 34
%% paperName: 2015年全国课标Ⅱ第 34 题 10 分
%% typeId: 6
%% type: 计算题
%% chapter: 机械波
%% point: 波的图象
%% method: 无
%% score: 10
%% degree: 450
%% duplicateId: 0
%% body:
【物理——选修 3-4】（15 分）

（1）（5 分）如图，一束光沿半径方向射向一块半圆柱形玻璃砖，在玻璃砖底面上的入射角为 $\theta$，经折射后射出 a、b 两束光线。则\xzanswer{ABD}（填正确答案标号。选对 1 个得 2 分，选对 2 个得 4 分，选对 3 个得 5 分。每选错 1 个扣 3 分，最低得分为 0 分）

\onepicture[h=2.8cm]{figs/34.png}

\fivechoices[answer=ABD]
{在玻璃中，a 光的传播速度小于 b 光的传播速度}
{在真空中，a 光的波长小于 b 光的波长}
{玻璃砖对 a 光的折射率小于对 b 光的折射率}
{若改变光束的入射方向使 $\theta$ 角变大，则折射光线 a 首先消失}
{分别用 a、b 光在同一个双缝干涉实验装置上做实验，a 光的干涉条纹间距大于 b 光的干涉条纹间距}

（2）（10 分）平衡位置位于原点 $O$ 的波源发出的简谐横波在均匀介质中沿水平 $x$ 轴传播，P、Q 为 $x$ 轴上的两个点（均位于 $x$ 轴正方向），P 与 $O$ 的距离为 $35\Ucm$，此距离介于一倍波长与二倍波长之间。已知波源自 $t=0$ 时由平衡位置开始向上振动，周期 $T=1\Us$，振幅 $A=5\Ucm$。当波传到 P 点时，波源恰好处于波峰位置；此后再经过 $5\Us$，平衡位置在 Q 处的质点第一次处于波峰位置。求：

\begin{enumerate}
	\item
	P、Q 间的距离；

	\item
	从 $t=0$ 开始到平衡位置在 Q 处的质点第一次处于波峰位置时，波源在振动过程中通过的路程。
\end{enumerate}

%% answer:
\jdanswer{
\begin{enumerate}
	\item
	$133\Ucm$

	\item
	$125\Ucm$
\end{enumerate}
}
%% memo:
\memoanswer{
	本题原卷及材料中未提供详解，待补充。
}

\item
%% number: 35
%% paperName: 2015年全国课标Ⅱ第 35 题 10 分
%% typeId: 6
%% type: 计算题
%% chapter: 运动和力
%% point: 动量和能量
%% method: 无
%% score: 10
%% degree: 450
%% duplicateId: 0
%% body:
【物理——选修 3-5】（15 分）

（1）（5 分）实物粒子和光都具有波粒二象性，下列事实中突出体现波动性的是\xzanswer{ACD}（填正确答案标号，选对 1 个给 2 分，选对 2 个得 4 分，选对 3 个得 5 分，每选错 1 个扣 3 分，最低得分 0 分）

\fivechoices[answer=ACD]
{电子束通过双缝实验后可以形成干涉图样}
{$\beta$ 射线在云室中穿过会留下清晰的径迹}
{人们利用慢中子衍射来研究晶体的结构}
{人们利用电子显微镜观测物质的微观结构}
{光电效应实验中，光电子的最大初动能与入射光的频率有关，与入射光的强度无关}

（2）（10 分）滑块 a、b 沿水平面上同一条直线运动，并发生碰撞；碰撞后两者粘在一起运动；经过一段时间后，从光滑路段进入粗糙路段。两者的位置 $x$ 随时间 $t$ 变化的图像如图所示。求：

\begin{enumerate}
	\item
	滑块 a、b 的质量之比；

	\item
	整个运动过程中，两滑块克服摩擦力做的功与因碰撞而损失的机械能之比。
\end{enumerate}

\onepicture[align=right, h=3.8cm]{figs/35.png}

%% answer:
\jdanswer{
\begin{enumerate}
	\item
	$\frac{1}{8}$

	\item
	$\frac{1}{2}$
\end{enumerate}
}
%% memo:
\memoanswer{
	本题原卷及材料中未提供详解，待补充。
}

\end{enumerate}

\end{document}
